\chapter{Metodologia}\label{chap:metodologia}
A metodologia deve permitir compreender e reproduzir o percurso científico. Recomenda-se explicitar hipóteses, formulação matemática, condições de contorno, discretização, algoritmos, parâmetros, infraestrutura e métricas de avaliação.

\section{Formulação do problema}
Considere um domínio $\Omega$ e uma variável de estado $\phi$. Um modelo genérico pode ser representado por
\begin{equation}\label{eq:modelo}\mathcal{L}(\phi;\boldsymbol{\theta})=f\quad\text{em }\Omega,\end{equation}
em que $\mathcal{L}$ é um operador, $\boldsymbol{\theta}$ representa parâmetros e $f$ é um termo de fonte. Esta equação é um exemplo abstrato, não uma descrição de um experimento realizado.

\section{Discretização e implementação}
O estudante deve identificar malha, passo temporal, tolerâncias, bibliotecas utilizadas e versões dos programas. Para modelos computacionais, registre critérios de convergência e a forma de verificação do código.

\section{Verificação e validação}
Verificação e validação respondem a perguntas distintas: se as equações discretizadas foram resolvidas corretamente e se o modelo representa adequadamente o fenômeno de interesse. Descreva testes, soluções de referência, incertezas e limitações.

\section{Reprodutibilidade}
Informe a origem dos dados, os procedimentos de processamento e a disponibilidade de código, quando pertinente. Use tabelas para reunir parâmetros e apresente claramente o ambiente computacional.
